\section*{Recorte e seleção dos artigos}

O tema do grupo é a desigualdade geográfica no acesso a modelos de inteligência artificial
generativa comercial. Para comparar artigos, estreitamos o tema para uma pergunta: como a
literatura empírica mede a desigualdade no acesso à IA e na adoção de tecnologias de fronteira,
entre organizações, entre empresas de um país e entre países.

A busca foi feita na OpenAlex, em título e resumo, em 11/09/2026, e está registrada no
Exercício~1. Pedro Dias Soares encontrou McElheran \emph{et al.} (2024), pela expressão \emph{AI
divide}, e Knez (2023), pela combinação \emph{cross-country}, \emph{technology adoption} e
\emph{diffusion}. Eduardo Velloso Landeiro encontrou Ahmed e Wahed (2020), pela expressão
\emph{compute divide}. O grupo é uma dupla, por isso Pedro ficou com dois artigos; cada um fichou o
que encontrou.

O conjunto segue a sugestão do enunciado. Ahmed e Wahed (2020) é o trabalho de referência: é o mais
antigo dos três e fixou o conceito de \emph{compute divide}. McElheran \emph{et al.} (2024) é o
trabalho recente, o primeiro retrato representativo da adoção de IA pelas empresas de um país
inteiro. Knez (2023) é a alternativa, porque explica a desigualdade por preços relativos
(salários), e não por acesso a insumos ou por características das organizações. Os três PDFs
acompanham este documento.
